\documentclass[a4paper, 12pt]{article}

% =========================================================
% PAQUETES
% =========================================================
\usepackage[table,xcdraw]{xcolor}
\usepackage[utf8]{inputenc}
\usepackage[T1]{fontenc}
\usepackage[spanish, es-nodecimaldot, es-tabla]{babel}
\usepackage{lmodern}
\usepackage{geometry}
\geometry{left=2.5cm, right=2.5cm, top=3cm, bottom=3cm}
\usepackage{graphicx}
\usepackage{booktabs}
\usepackage{longtable}
\usepackage{array}
\usepackage{float}
\usepackage{hyperref}
\usepackage{fancyhdr}
\usepackage{titlesec}
\usepackage{enumitem}
\usepackage{listings}
\usepackage{tikz}
\usepackage{amsmath}
\usepackage{tabularx}
\usepackage{setspace}

% =========================================================
% CONFIGURACIÓN DE PÁGINA Y ENLACES
% =========================================================
\definecolor{ucsblue}{HTML}{004B87}
\definecolor{codegray}{rgb}{0.95,0.95,0.95}

\hypersetup{
    colorlinks=true,
    linkcolor=ucsblue,
    urlcolor=ucsblue,
    citecolor=ucsblue
}

\pagestyle{fancy}
\fancyhf{}
\rhead{SDD - Sistema de Gestión de Droguería}
\lhead{Documento de Diseño de Software}
\cfoot{\thepage}

% =========================================================
% CONFIGURACIÓN DE CÓDIGO
% =========================================================
\lstset{
    backgroundcolor=\color{codegray},
    basicstyle=\ttfamily\small,
    breaklines=true,
    frame=single,
    keywordstyle=\color{blue},
    commentstyle=\color{green!50!black},
    stringstyle=\color{red!70!black},
    numbers=none,
    captionpos=b,
    tabsize=2,
    showstringspaces=false,
    inputencoding=utf8,
    extendedchars=true,
    literate={á}{{\'a}}1 {é}{{\'e}}1 {í}{{\'i}}1 {ó}{{\'o}}1 {ú}{{\'u}}1 {ñ}{{\~n}}1
}

\begin{document}

% =========================================================
% PORTADA (Exactamente el formato solicitado)
% =========================================================
\begin{titlepage}
    \begin{flushleft}
        % Asegúrate de tener el archivo logo.png en la misma carpeta
        \includegraphics[width=3cm]{logo.png} 
    \end{flushleft}
    \centering
    
    \vspace{1cm}
    {\Large \textbf{Universidad Católica de la Santísima Concepción}}\\[3cm]
    
    {\Huge \textbf{Ingeniería de Ejecución Informática}}\\[0.5cm]
    {\Large \textbf{Documento de Diseño de Software (SDD)\\[0.2cm] Plataforma de Gestión Inteligente de la Droguería Comunal}}\\[0.5cm]
    
    \vfill
    
    \begin{minipage}{0.45\textwidth}
        \raggedright
        \textbf{Supervisor Institucional:}\\
        Pablo Cáceres Sáez \\
        \textbf{Profesor Evaluador:}\\
        Nicolás Pérez \\
    \end{minipage}
    \hfill
    \begin{minipage}{0.45\textwidth}
        \raggedleft
        \textbf{Autor / Desarrollador:}\\
        Cristóbal Matías Faúndez Stuardo \\
        \textbf{Organización:}\\
        DAS Talcahuano \\
    \end{minipage}
    
    \vfill
    {\large 07 de octubre de 2026}
\end{titlepage}

% =========================================================
% CONTROL DE VERSIONES (Solucionado el desbordamiento)
% =========================================================
\newpage
\section*{Control de versiones}

\vspace{0.5cm}
\noindent
\renewcommand{\arraystretch}{1.5} % Da un poco más de espacio vertical a las celdas
\begin{tabularx}{\textwidth}{@{} l l >{\raggedright\arraybackslash}p{4.5cm} >{\raggedright\arraybackslash}X @{}}
\toprule
\textbf{Versión} & \textbf{Fecha} & \textbf{Autor} & \textbf{Descripción} \\
\midrule
0.1 & 25 de Ago 2026 & Cristóbal Matías Faúndez Stuardo & Creación de la estructura base del documento. \\
0.5 & 05 de Sep 2026 & Cristóbal Matías Faúndez Stuardo & Documentación de los módulos Auth y Dashboard. \\
0.8 & 18 de Sep 2026 & Cristóbal Matías Faúndez Stuardo & Incorporación de vistas de Inventario y Despachos. \\
0.9 & 30 de Sep 2026 & Cristóbal Matías Faúndez Stuardo & Mapeo de errores y módulo de Trazabilidad. \\
1.0 & 07 de Oct 2026 & Cristóbal Matías Faúndez Stuardo & Revisión final, correcciones ortográficas y versión aprobada. \\
\bottomrule
\end{tabularx}
\renewcommand{\arraystretch}{1} % Restaura el espacio vertical normal

\newpage
% =========================================================
% ÍNDICE
% =========================================================
\tableofcontents
\newpage

% =========================================================
% 1. INTRODUCCIÓN
% =========================================================

\section{Introducción}

El presente Documento de Diseño de Software (SDD) describe la arquitectura, los módulos, la base de datos y las reglas de negocio de la \textbf{Plataforma de Gestión Inteligente de la Droguería Comunal} desarrollada para la Dirección de Administración de Salud (DAS) de la Ilustre Municipalidad de Talcahuano. El sistema digitaliza el ciclo completo de abastecimiento farmacéutico comunal: \textbf{recepción de mercadería, bodega de inventario por lotes, cuarentena sanitaria, despachos normales y de estupefacientes, solicitudes de establecimientos, distribución física en bodega, trazabilidad por lote y reportes valorizados}.

La plataforma reemplaza planillas dispersas y registros manuales por un sistema transaccional con metodología \textbf{FEFO} (\textit{First Expired, First Out}), trazabilidad inmutable de movimientos, control legal de estupefacientes y control de acceso por roles y permisos por vista.

% =========================================================
\section{Propósito del documento}

Este documento está dirigido al equipo técnico responsable del mantenimiento y evolución de la plataforma. Su propósito es servir como fuente única de verdad sobre:
\begin{itemize}
    \item La arquitectura de software y las tecnologías empleadas.
    \item El modelo de datos y las relaciones entre tablas.
    \item Las reglas de negocio críticas (FEFO, estados de lote, restricciones de estupefacientes).
    \item La estructura de roles y permisos.
    \item Los puntos de integración (API interna, exportaciones, generación de PDF y códigos QR).
\end{itemize}

% =========================================================
\section{Objetivos}

\subsection{Objetivo general}
Describir de manera completa y precisa el diseño de software de la plataforma de gestión de la droguería comunal DAS Talcahuano, de modo que cualquier desarrollador pueda mantener, extender y auditar el sistema sin conocimiento previo del código.

\subsection{Objetivos específicos}
\begin{enumerate}
    \item Documentar la arquitectura en capas y el flujo transaccional de datos.
    \item Especificar el modelo de base de datos con todas sus tablas y relaciones.
    \item Definir las reglas de negocio: FEFO, estados de lote (Disponible, Cuarentena, Merma), segregación de despachos normales y controlados.
    \item Documentar el sistema de roles y permisos por vista basado en Spatie Permission.
    \item Establecer los estándares de seguridad, respaldo y mantenimiento.
\end{enumerate}

% =========================================================
\section{Alcance}

\subsection{Incluido en el proyecto}
\begin{itemize}
    \item Autenticación con recuperación de contraseña administrada.
    \item Mantenedor de proveedores.
    \item Recepción de mercadería con generación de PDF y registro en inventario por lote.
    \item Inventario valorizado con carga masiva Excel, filtros, código QR por lote y modales de gestión de estado.
    \item Cuarentena sanitaria y mermas con motivo y responsable.
    \item Despacho normal con FEFO server-side y descuento de stock transaccional.
    \item Despacho controlado exclusivo para estupefacientes con series de recetas cheque.
    \item Solicitudes de establecimientos con ruteo automático normal/controlado según tipo de producto.
    \item Visor 3D de distribución con capacidad, niveles, posiciones y ocupación porcentual.
    \item Trazabilidad consolidada sobre la tabla \texttt{stock\_movements}.
    \item Reportes valorizados exportables a Excel y PDF.
    \item Administración de usuarios, roles y accesos por vista.
    \item Notificaciones clínicas (vencimientos, cuarentena, stock bajo) y búsqueda global.
\end{itemize}

\subsection{Fuera del alcance}
\begin{itemize}
    \item Prescripción médica electrónica y receta clínica digital.
    \item Facturación electrónica al SII (la recepción registra el documento de respaldo, no lo valida).
    \item Integración automática con el sistema de compras municipales.
    \item Aplicación móvil nativa (la interfaz es web responsiva).
    \item Control de cadenas de frío por sensores IoT.
\end{itemize}

% =========================================================
\section{Descripción general del sistema}

La plataforma es una aplicación web \textbf{Laravel} (patrón MVC) con base de datos \textbf{MySQL} y front-end \textbf{Blade + Tailwind CSS + Alpine.js}. El flujo de negocio central es:

\begin{enumerate}
    \item \textbf{Recepción}: se registra la mercadería entrante (documento, proveedor, lotes, vencimientos, valores). Cada lote se crea o acumula en \texttt{inventories} en estado \textbf{Cuarentena} y se registra un movimiento en \texttt{stock\_movements}.
    \item \textbf{Liberación}: el personal autorizado cambia el estado del lote a \textbf{Disponible} (o Merma) desde Inventario o Cuarentena, con motivo documentado.
    \item \textbf{Solicitudes}: los establecimientos (CESFAM y otros centros) registran pedidos. Solo los CESFAM pueden incluir estupefacientes.
    \item \textbf{Despacho}: los medicamentos/insumos salen por el módulo de Despacho Normal; los estupefacientes por el módulo de Despacho Controlado. Ambos descuentan inventario con FEFO, generan guía PDF y registran trazabilidad.
    \item \textbf{Distribución}: el visor 3D refleja en tiempo real el stock por estantería, nivel y posición, con porcentaje de ocupación sobre la capacidad definida.
    \item \textbf{Trazabilidad y Reportes}: toda operación queda registrada en \texttt{stock\_movements} con usuario, documento y lote; los reportes consolidan y valorizan el stock disponible.
\end{enumerate}

% =========================================================
\section{Requisitos funcionales}

\begin{longtable}{@{}p{1.2cm} p{4.2cm} X@{}}
\toprule
\textbf{Código} & \textbf{Módulo} & \textbf{Descripción} \\
\midrule
\endhead
RF-01 & Autenticación & Login por correo y contraseña con regeneración de sesión. Solicitud de restablecimiento desde el login (throttled) y asignación de nueva clave por el administrador. \\
RF-02 & Autenticación & Recuperación con estado \texttt{pending}/\texttt{resolved} en \texttt{password\_reset\_tokens}. \\
RF-03 & Proveedores & CRUD con RUT único, filtros por búsqueda, categoría y estado, KPIs de actividad y activación/desactivación. \\
RF-04 & Recepción & Registro de cabecera (documento, proveedor, origen) y detalle dinámico de productos (lote, vencimiento, cantidad, valor, categoría, programa). Correlativo automático con bloqueo transaccional. Generación de PDF. \\
RF-05 & Recepción & Cada lote se crea o acumula en inventario según (producto, lote) y se registra movimiento de trazabilidad con usuario. \\
RF-06 & Inventario & Agrupación por producto con desglose por lote (ubicación, estado, valor unitario, valor total, programa, nivel, posición). \\
RF-07 & Inventario & Carga masiva desde plantilla Excel con validación por fila, acumulación de lotes y enlaces automáticos a productos. \\
RF-08 & Inventario & Cambio de estado del lote (Disponible/Cuarentena/Merma) con motivo obligatorio en bloqueos. \\
RF-09 & Inventario & Etiqueta QR por lote con contenido estructurado escaneable por el módulo de despacho. \\
RF-10 & Cuarentena & Vista dedicada a lotes retenidos con liberación, merma o mantención, y registro histórico de mermas valorizado. \\
RF-11 & Solicitudes & Registro de solicitudes con establecimiento validado, responsable, firma digital del solicitante y productos según el tipo de establecimiento. \\
RF-12 & Solicitudes & Solo establecimientos CESFAM pueden solicitar estupefacientes (validación cliente y servidor). \\
RF-13 & Despacho normal & Selección FEFO de lotes, validación de stock, descuento transaccional, guía de despacho PDF y ruteo de estupefacientes al módulo controlado. \\
RF-14 & Despacho controlado & Exclusivo para la lista de fármacos controlados; acepta solo solicitudes CESFAM; registra series de recetas cheque; firma del solicitante precargada. \\
RF-15 & Solicitudes & Al usar una solicitud, \texttt{cantidad\_despachada} se actualiza y el estado evoluciona Pendiente $\to$ Parcial $\to$ Despachado. \\
RF-16 & Distribución & Visor 3D con pasillos coloreados, estanterías con postes/repisas, niveles y posiciones reales de la BD, capacidad por estantería y porcentaje de ocupación en vivo. \\
RF-17 & Distribución & Arrastrar estanterías actualiza sus coordenadas (pos\_x, pos\_z) vía AJAX. \\
RF-18 & Trazabilidad & Listado consolidado de movimientos (recepciones y despachos) con filtros, paginación y exportación; línea de tiempo por lote. \\
RF-19 & Reportes & Reporte valorizado del stock disponible con filtros (búsqueda, categoría, pasillo, estante, programa) y exportación Excel/PDF con totales. \\
RF-20 & Usuarios & CRUD de usuarios con roles; administración de roles con accesos por vista; solicitudes de recuperación pendientes. \\
RF-21 & Seguridad & Acceso por permisos: cada módulo exige \texttt{ver.<modulo>}; el rol Administrador posee acceso total. \\
RF-22 & Notificaciones & Alertas de vencimientos (60 días), lotes en cuarentena y stock bajo disponible (cantidad $\le$ mínimo). \\
RF-23 & Búsqueda global & Búsqueda unificada de stock, catálogo y proveedores con atajo Ctrl+K. \\
RF-24 & Dashboard & KPIs de valor de inventario, unidades, recepciones/despachos del período, próximos a vencer, cuarentena y últimos movimientos. \\
\bottomrule
\end{longtable}

% =========================================================
\section{Requisitos no funcionales}

\begin{itemize}
    \item \textbf{Rendimiento}: paginación server-side en listados; consultas agrupadas en SQL; debouncing de 300 ms y AbortController en la búsqueda global.
    \item \textbf{Integridad}: transacciones con \texttt{DB::transaction} y bloqueos \texttt{lockForUpdate} en operaciones de stock; claves foráneas con \texttt{onDelete} definido.
    \item \textbf{Seguridad}: hashing bcrypt, validación server-side en todos los escritos, tokens CSRF en todos los formularios, \texttt{throttle} en rutas públicas y consultas Eloquent parametrizadas (sin inyección SQL).
    \item \textbf{Usabilidad}: interfaz responsiva (móvil/escritorio), patrón visual unificado (migas de pan, tarjetas, modales), validaciones en línea.
    \item \textbf{Auditabilidad}: todo movimiento de stock registra usuario y documento de respaldo.
    \item \textbf{Mantenibilidad}: lista de estupefacientes y establecimientos centralizadas en \texttt{config/drogueria.php}; servicios desacoplados (\texttt{InventoryService}, \texttt{SolicitudService}).
    \item \textbf{Compatibilidad}: navegadores modernos (Chrome, Edge, Firefox); cámara QR requiere HTTPS o localhost.
\end{itemize}

% =========================================================
\section{Tecnologías utilizadas}

\subsection{Resumen tecnológico}
\begin{table}[H]
\centering
\begin{tabularx}{\textwidth}{|X|X|X|}
\hline
\rowcolor{ucsblue} \textcolor{white}{\textbf{Capa}} & \textcolor{white}{\textbf{Tecnología}} & \textcolor{white}{\textbf{Uso}} \\ \hline
Back-end & PHP 8.3+ / Laravel 13 & Controladores, servicios, validación, migraciones, middleware. \\ \hline
Base de datos & MySQL 8 & Persistencia transaccional con claves foráneas. \\ \hline
Front-end & Blade, Tailwind CSS 3.4, Alpine.js & Vistas, estilos utilitarios, estado reactivo ligero. \\ \hline
Librerías JS & Three.js, SignaturePad, SweetAlert2, html5-qrcode, qrcodejs & Visor 3D, firmas digitales, alertas, escáner QR. \\ \hline
Paquetes PHP & barryvdh/laravel-dompdf, maatwebsite/excel, simplesoftwareio/simple-qrcode, spatie/laravel-permission & PDFs, importación/exportación Excel, generación QR, roles y permisos. \\ \hline
\end{tabularx}
\caption{Stack tecnológico de la plataforma.}
\end{table}

\subsection{Laravel 13}
Núcleo del sistema. Se utilizan migraciones versionadas, Eloquent ORM con relaciones y castings, Form Validation, Blade como motor de plantillas con layout principal (\texttt{layouts/app.blade.php}), middleware de autenticación, rol y permiso, y el contenedor de servicios para registrar el bypass de permisos del rol Administrador (\texttt{Gate::before}).

\subsection{Tailwind CSS 3.4 y Alpine.js}
Estilos utilitarios con paleta institucional (teal/slate) y componentes reutilizables. Alpine.js gobierna la interactividad del visor 3D (paneles, tooltips, cámara). El plugin \texttt{@tailwindcss/forms} unifica formularios.

\subsection{MySQL 8}
Motor transaccional InnoDB con claves foráneas, índices compuestos para búsqueda (producto, lote, estado, vencimiento) y columnas calculadas de valorización resueltas en consulta (\texttt{SUM(cantidad * valor\_neto\_unitario)}).

\subsection{Librerías de interfaz}
\begin{itemize}
    \item \textbf{Three.js r128 + OrbitControls/DragControls}: visor 3D de la bodega con modo 2D/3D, zoom, pantalla completa y arrastre de estanterías.
    \item \textbf{SignaturePad 4}: captura de firmas (despacha y solicita/recibe) en canvas, con precarga de la firma del solicitante mediante \texttt{fromDataURL}.
    \item \textbf{SweetAlert2}: confirmaciones y avisos operativos.
    \item \textbf{html5-qrcode}: escaneo de etiquetas QR en despacho.
\end{itemize}

% =========================================================
\section{Gestión del proyecto}

\subsection{Trello}
El flujo de trabajo se gestiona en tablero Kanban (Backlog, En progreso, Testing, Terminado) con tarjetas por módulo.

\subsection{GitHub}
Repositorio \texttt{Proyecto\_Drogueria} bajo control de versiones Git. Las funcionalidades se integran por ramas por módulo con revisiones locales previas al merge.

% =========================================================
\section{Arquitectura del sistema}

\subsection{Capa de presentación}
Vistas Blade organizadas por módulo en \texttt{resources/views}: \texttt{auth}, \texttt{dashboard}, \texttt{reception}, \texttt{inventario}, \texttt{quarantine}, \texttt{dispatch}, \texttt{controlled\_dispatch}, \texttt{requests}, \texttt{distribution} (con mantenedores), \texttt{traceability}, \texttt{reportes}, \texttt{users}, \texttt{providers}, \texttt{profile}, \texttt{errors}. El layout principal provee sidebar con filtrado por permisos, notificaciones, búsqueda global y toasts de flash messages.

\subsection{Capa de aplicación}
\begin{itemize}
    \item \textbf{Controladores HTTP}: un controlador por módulo en \texttt{app/Http/Controllers} (18 controladores, incluidos los de autenticación Breeze personalizados).
    \item \textbf{Servicios de dominio}:
    \begin{itemize}
        \item \texttt{InventoryService}: consumo de stock con FEFO y \texttt{lockForUpdate}, ingreso/acumulación de lotes, registro de movimientos, helpers de controlados y establecimientos.
        \item \texttt{SolicitudService}: registro de despachos contra solicitudes y recálculo de estado global (Pendiente/Parcial/Despachado).
    \end{itemize}
    \item \textbf{Importación}: \texttt{App\textbackslash Imports\textbackslash InventoriesImport} (Maatwebsite Excel) con validación fila a fila.
    \item \textbf{Exportación}: \texttt{App\textbackslash Exports\textbackslash InventoryExport}.
\end{itemize}

\subsection{Capa de datos}
MySQL 8 con 25+ tablas administradas por migraciones. Modelos Eloquent con \texttt{\$fillable} explícito, castings y relaciones declaradas.

\subsection{Procesamiento auxiliar}
Generación de PDF (DomPDF) para guías y reportes; importación/exportación Excel; generación de códigos QR en SVG; toasts de flash messages por sesión.

% =========================================================
\section{Diseño de módulos}

\subsection{Módulo de autenticación}
Login propio (\texttt{LoginController}) con regeneración de sesión y bloqueo de intentos vía throttle. Las solicitudes de recuperación se registran en \texttt{password\_reset\_tokens} con estado \texttt{pending}; el administrador asigna la nueva clave desde \texttt{Users}, quedando \texttt{resolved}. El formulario de login incluye mostrar/ocultar contraseña y validaciones de formato y largo (8--64).

\subsection{Módulo de proveedores}
Mantenedor con RUT único, búsqueda, filtros por categoría/estado/mes, KPIs (total, activos, agregados del mes), alta por modal y cambio de estado activo/inactivo por AJAX.

\subsection{Módulo de recepción}
Formulario de dos secciones: datos del documento (tipo, número, orden de compra, origen, proveedor, firma del receptor) y detalle dinámico de productos (producto del catálogo, categoría operacional, programa, lote, vencimiento, cantidad, valor neto). Al guardar, dentro de una transacción: correlativo automático con bloqueo, creación de la recepción, items, \textbf{ingreso a inventario por lote} (crea o acumula) en estado Cuarentena, movimiento de trazabilidad y cálculo del gran total con IVA 19\% (configurable). Descarga de PDF de la recepción.

\subsection{Módulo de inventario}
Tabla agrupada por producto con desglose por lote: ubicación (pasillo/estante/nivel/posición), chips de programa y flags de pedido, estado sanitario, stock y \textbf{valorización} (valor unitario y total por lote, total por producto). Filtros: búsqueda, categoría, pasillo, estante y programa. Acciones por lote: cambio de estado (modal valorizado), etiqueta QR (generada server-side en SVG). Ingreso directo manual y \textbf{carga masiva Excel} con plantilla descargable (19 columnas validadas). Modal de mermas valorizado con total.

\subsection{Módulo de cuarentena}
Listado de lotes en estado Cuarentena con búsqueda, liberación, declaración de Merma o mantención en cuarentena (modal valorizado) y tabla de mermas históricas con total valorizado.

\subsection{Módulo de despacho normal}
Selección FEFO: el listado del formulario proviene del inventario en estado Disponible (excluyendo controlados) ordenado por vencimiento ascendente; el cliente preselecciona el lote más próximo a vencer y el servidor \textbf{valida y consume con FEFO} (reparto entre lotes si es necesario) dentro de transacción con bloqueo. Incluye clasificación interno/externo, condiciones de transporte, firmas digitales, escáner QR de etiquetas, panel de solicitudes pendientes (los estupefacientes se excluyen y se avisa al usuario) y bandeja histórica con PDF.

\subsection{Módulo de despacho controlado}
Restringido a los fármacos de la lista de controlados provenientes del inventario en estado Disponible. Correlativo GDC-AAAA-NNNN, series de recetas cheque, firma de quien recibe precargada desde la solicitud CESFAM, y modal propio de solicitudes con estupefacientes pendientes. El descuento de stock y la trazabilidad (\texttt{Despacho Controlado}) siguen la misma mecánica FEFO.

\subsection{Módulo de solicitudes}
Registro por establecimiento (lista oficial agrupada CESFAM/Otros), responsable, cargo, observaciones, firma digital del solicitante y detalle de productos con consumo anterior y stock actual. Restricción de estupefacientes solo para CESFAM (validación de cliente y servidor). Ruteo automático: al usar la solicitud en despacho normal solo se cargan productos no controlados; en despacho controlado solo estupefacientes. Estado global recalculado por \texttt{SolicitudService}.

\subsection{Módulo de distribución}
Visor 3D (Three.js) que representa la bodega real: zonas de pasillo coloreadas, estanterías con postes, repisas, \textbf{niveles numerados} y \textbf{posiciones nombradas} desde la base de datos, slots coloreados según ocupación real del inventario, capacidad configurable por estantería y \textbf{porcentaje de ocupación} en marcadores flotantes, leyenda y tooltip con barra de progreso. Mantenedores de pasillos, estanterías (con capacidad), niveles y posiciones. Botón ``Ver detalle'' redirige al inventario filtrado por estantería.

\subsection{Módulo de trazabilidad}
Fuente única \texttt{stock\_movements}: listado paginado con filtros (búsqueda, tipo de movimiento, rango de fechas), KPIs del mes, exportación Excel y línea de tiempo por lote (JSON) con ficha completa del lote (proveedor, documento, valor, ubicación, estado, stock).

\subsection{Módulo de reportes}
Reporte valorizado del stock Disponible: KPI de valor total según filtros, tabla agrupada por producto con desglose por lote valorizado y exportaciones Excel y PDF (con fila de total general).

\subsection{Módulo de usuarios y roles}
CRUD de usuarios (con apellido y rol) y administración de \textbf{roles con accesos por vista}: 12 módulos configuran permisos \texttt{ver.<modulo>} sincronizados vía Spatie. El rol Administrador posee acceso total (\texttt{Gate::before}). Las solicitudes de recuperación pendientes se gestionan desde este módulo.

\subsection{Módulo de notificaciones y búsqueda global}
Campana con conteo y listado de alertas clínicas (vencimientos a 60 días, cuarentena, stock bajo), endpoint JSON \texttt{/api/notifications}; búsqueda global \texttt{/api/global-search} sobre inventario, catálogo y proveedores con resaltado y navegación por teclado (Ctrl+K).

% =========================================================
\section{Diseño de la base de datos}

El modelo de datos se organiza en cinco dominios: seguridad, catálogo, abastecimiento, salidas y distribución física.

\subsection{Tabla users}
\begin{table}[H]
\centering
\begin{tabularx}{\textwidth}{|l|l|X|}
\hline
\rowcolor{ucsblue} \textcolor{white}{\textbf{Columna}} & \textcolor{white}{\textbf{Tipo}} & \textcolor{white}{\textbf{Descripción}} \\ \hline
id & BIGINT PK & Identificador. \\ \hline
name / last\_name & VARCHAR & Nombres y apellidos del funcionario. \\ \hline
email & VARCHAR UNIQUE & Credencial de acceso. \\ \hline
password & VARCHAR & Hash bcrypt (cast \texttt{hashed}). \\ \hline
email\_verified\_at & TIMESTAMP & Verificación opcional. \\ \hline
\end{tabularx}
\caption{Tabla \texttt{users}.}
\end{table}

\subsection{Tabla products (catálogo maestro)}
\begin{table}[H]
\centering
\begin{tabularx}{\textwidth}{|l|l|X|}
\hline
\rowcolor{ucsblue} \textcolor{white}{\textbf{Columna}} & \textcolor{white}{\textbf{Tipo}} & \textcolor{white}{\textbf{Descripción}} \\ \hline
id & BIGINT PK & Identificador. \\ \hline
nombre & VARCHAR & Nombre oficial del medicamento o insumo. \\ \hline
\end{tabularx}
\caption{Tabla \texttt{products}. La recepción y las solicitudes seleccionan desde este catálogo; los despachos seleccionan desde \texttt{inventories} (stock real).}
\end{table}

\subsection{Tabla providers}
\texttt{id}, \texttt{rut} (único), \texttt{name} (razón social), \texttt{business\_line}, \texttt{category}, \texttt{email}, \texttt{phone}, \texttt{address}, \texttt{status} (\texttt{activo}/\texttt{inactivo}), timestamps.

\subsection{Tablas receptions y reception\_items}
\textbf{receptions}: \texttt{id}, \texttt{correlativo} (único, autogenerado), \texttt{fecha\_recepcion}, \texttt{tipo\_documento}, \texttt{numero\_documento}, \texttt{orden\_compra}, \texttt{origen}, \texttt{provider\_id} (FK), \texttt{valor\_total\_recepcion} (IVA incluido), \texttt{firma\_proveedor}, \texttt{firma\_receptor}, timestamps.

\textbf{reception\_items}: \texttt{id}, \texttt{reception\_id} (FK cascada), \texttt{product\_id} (FK), \texttt{tipo\_producto}, \texttt{lote}, \texttt{fecha\_vencimiento}, \texttt{cantidad}, \texttt{valor\_unitario}, \texttt{valor\_total\_item} (IVA incluido).

\subsection{Tabla inventories (núcleo del sistema)}
Un registro por \textbf{producto + lote} con el stock físico actual.
\begin{table}[H]
\centering
\small
\begin{tabularx}{\textwidth}{|l|l|X|}
\hline
\rowcolor{ucsblue} \textcolor{white}{\textbf{Columna}} & \textcolor{white}{\textbf{Tipo}} & \textcolor{white}{\textbf{Descripción}} \\ \hline
id & BIGINT PK & Identificador del lote. \\ \hline
product\_id & BIGINT FK & Catálogo maestro (normalizado por migración). \\ \hline
user\_id & BIGINT FK & Último usuario que ingresó o modificó el lote. \\ \hline
correlativo\_interno & INT & Correlativo del documento de recepción. \\ \hline
fecha\_recepcion / fecha\_vencimiento & DATE & Fechas sanitarias. \\ \hline
tipo\_documento / numero\_documento / orden\_compra & VARCHAR & Respaldo documental. \\ \hline
origen\_mercaderia / rut\_proveedor / razon\_social & VARCHAR & Procedencia (desnormalizada para trazabilidad). \\ \hline
producto / categoria & VARCHAR & Nombre y categoría operacional (F-FOFAR, I-URGENCIA, etc.). \\ \hline
estado & VARCHAR & \texttt{Disponible}, \texttt{Cuarentena} o \texttt{Merma}. \\ \hline
motivo & TEXT & Razón documentada en cuarentena/merma. \\ \hline
lote & VARCHAR & Código de lote (\texttt{S/L} si no trae). \\ \hline
cantidad / minimo & INT & Stock físico y umbral mínimo para alerta de reposición. \\ \hline
valor\_neto\_unitario & DECIMAL(15,2) & Valor neto por unidad (valorización). \\ \hline
pasillo / estante / nivel / posicion & VARCHAR & Ubicación física (enlaza con la distribución 3D por nombre). \\ \hline
programa & VARCHAR & Programa de salud (Morbilidad General, Cardiovascular, IRA/ERA). \\ \hline
solicitud\_pedido\_mensual, solicitud\_stock\_critico, adelanto\_de\_pedido, traspaso & BOOLEAN & Marcas de gestión de pedidos. \\ \hline
\end{tabularx}
\caption{Tabla \texttt{inventories}. Índices en \texttt{producto}, \texttt{lote}, \texttt{estado} y \texttt{fecha\_vencimiento}.}
\end{table}

\subsection{Tablas dispatches y dispatch\_items}
\textbf{dispatches}: \texttt{id}, \texttt{correlativo} (único, GD-AAAA-NNNN), \texttt{fecha\_despacho}, \texttt{clasificacion} (Interno/Externo), \texttt{estado\_liberacion}, \texttt{establecimiento\_destino}, \texttt{direccion\_destino}, \texttt{condiciones\_transporte}, \texttt{firma\_despacha}, \texttt{firma\_recibe}.

\textbf{dispatch\_items}: \texttt{id}, \texttt{dispatch\_id} (FK), \texttt{inventory\_id} (FK al lote consumido), \texttt{producto}, \texttt{lote}, \texttt{fecha\_vencimiento}, \texttt{programa}, \texttt{cantidad}.

\subsection{Tablas controlled\_dispatches y controlled\_dispatch\_items}
Idénticas a las anteriores con \texttt{correlativo} GDC-AAAA-NNNN y \texttt{series\_recetas} (talonario de recetas cheque).

\subsection{Tablas establishment\_requests y establishment\_items}
\textbf{establishment\_requests}: \texttt{id}, \texttt{establecimiento}, \texttt{fecha\_solicitud}, \texttt{numero\_olicitud} (\#NNNNN), \texttt{responsable}, \texttt{cargo}, \texttt{observaciones}, \texttt{firma\_solicitante} (PNG base64), \texttt{estado} (Pendiente/Parcial/Despachado).

\textbf{establishment\_items}: \texttt{id}, \texttt{establishment\_request\_id} (FK cascada), \texttt{product\_id} (FK), \texttt{consumo\_anterior}, \texttt{stock\_actual}, \texttt{cantidad\_solicitada}, \texttt{cantidad\_despachada}.

\subsection{Tabla stock\_movements (trazabilidad)}
\texttt{id}, \texttt{inventory\_id} (FK), \texttt{user\_id} (FK), \texttt{producto}, \texttt{lote}, \texttt{fecha\_vencimiento}, \texttt{tipo\_movimiento} (Recepción / Despacho / Despacho Controlado), \texttt{cantidad} (positiva/negativa), \texttt{fecha\_hora}, \texttt{origen\_destino}, \texttt{documento\_respaldo}, \texttt{estado\_lote}. Índice compuesto (\texttt{producto}, \texttt{lote}).

\subsection{Tablas de distribución física}
\begin{itemize}
    \item \textbf{aisles}: \texttt{id}, \texttt{name}, \texttt{description}, \texttt{color} (hex, color del pasillo en el 3D).
    \item \textbf{shelves}: \texttt{id}, \texttt{name} (coincide con \texttt{inventories.estante}), \texttt{aisle\_id} (FK), \texttt{family\_category}, \texttt{pos\_x}, \texttt{pos\_z}, \texttt{rotation}, \texttt{capacidad} (unidades, 0 = sin límite).
    \item \textbf{levels}: \texttt{id}, \texttt{shelf\_id} (FK), \texttt{number}, \texttt{description}.
    \item \textbf{positions}: \texttt{id}, \texttt{level\_id} (FK), \texttt{name}, \texttt{description}.
\end{itemize}

\subsection{Tablas de seguridad (Spatie Permission)}
\texttt{permissions} (\texttt{name} = \texttt{ver.<módulo>}), \texttt{roles}, \texttt{model\_has\_roles}, \texttt{model\_has\_permissions}, \texttt{role\_has\_permissions}, además de \texttt{password\_reset\_tokens} (con columna \texttt{status}) y \texttt{sessions}.

% =========================================================
\section{Relaciones principales}
\begin{itemize}
    \item \texttt{receptions.provider\_id} $\to$ \texttt{providers.id} (restrict).
    \item \texttt{reception\_items.reception\_id} $\to$ \texttt{receptions.id} (cascade) y \texttt{reception\_items.product\_id} $\to$ \texttt{products.id}.
    \item \texttt{inventories.product\_id} $\to$ \texttt{products.id} (set null) y \texttt{inventories.user\_id} $\to$ \texttt{users.id}.
    \item \texttt{dispatch\_items.dispatch\_id} $\to$ \texttt{dispatches.id} (cascade) y \texttt{dispatch\_items.inventory\_id} $\to$ \texttt{inventories.id} (set null) --- enlaza cada ítem despachado con el lote físico consumido.
    \item \texttt{controlled\_dispatch\_items.controlled\_dispatch\_id} $\to$ \texttt{controlled\_dispatches.id} (cascade) y \texttt{inventory\_id} $\to$ \texttt{inventories.id}.
    \item \texttt{establishment\_items.establishment\_request\_id} $\to$ \texttt{establishment\_requests.id} (cascade) y \texttt{product\_id} $\to$ \texttt{products.id}.
    \item \texttt{stock\_movements.inventory\_id} $\to$ \texttt{inventories.id} y \texttt{user\_id} $\to$ \texttt{users.id}.
    \item \texttt{shelves.aisle\_id} $\to$ \texttt{aisles.id}; \texttt{levels.shelf\_id} $\to$ \texttt{shelves.id}; \texttt{positions.level\_id} $\to$ \texttt{levels.id}.
\end{itemize}

% =========================================================
\section{Casos de uso}

\subsection{Actores}
\begin{itemize}
    \item \textbf{Bodeguero / Operador}: recepciona, gestiona inventario y cuarentena, despacha.
    \item \textbf{Químico Farmacéutico}: libera o manda a merma lotes, autoriza despachos controlados.
    \item \textbf{Administrador (DAS)}: acceso total, usuarios, roles, proveedores y configuración.
    \item \textbf{Establecimiento (CESFAM/otros)}: actor externo que genera solicitudes (registradas por el operador).
\end{itemize}

\subsection{Casos de uso principales}
\begin{enumerate}
    \item Iniciar sesión / solicitar recuperación.
    \item Registrar proveedor.
    \item Recepcionar mercadería con lotes.
    \item Importar inventario masivamente.
    \item Liberar cuarentena / declarar merma.
    \item Registrar solicitud de establecimiento.
    \item Despachar con FEFO (normal).
    \item Despachar estupefacientes (controlado).
    \item Generar etiqueta QR / escanear lote.
    \item Consultar trazabilidad por lote.
    \item Generar reportes valorizados.
    \item Administrar usuarios, roles y accesos.
    \item Configurar distribución física (pasillos, estanterías, capacidad).
\end{enumerate}

% =========================================================
\section{Caso de uso detallado: Despachar con FEFO}

\subsection{Flujo principal}
\begin{enumerate}
    \item El operador abre \textbf{Salida $\to$ Despacho Normal}. El sistema muestra el inventario Disponible ordenado por vencimiento (FEFO).
    \item Selecciona clasificación, destino, fecha y liberación; agrega productos (el lote FEFO se autocompleta) o carga una solicitud.
    \item Confirma. El servidor abre transacción, genera correlativo con bloqueo y, por cada producto, ejecuta \texttt{InventoryService::consumir()}:
    \begin{itemize}
        \item Bloquea los lotes candidatos (\texttt{lockForUpdate}) del producto en estado Disponible.
        \item Ordena por \texttt{fecha\_vencimiento} ascendente (nulos al final) y descuenta desde el más próximo a vencer, repartiendo entre lotes si la cantidad excede al primero.
        \item Si el stock total es insuficiente, aborta con mensaje explícito y hace \texttt{rollback}.
    \end{itemize}
    \item Por cada asignación crea un \texttt{dispatch\_item} (con \texttt{inventory\_id}) y un \texttt{stock\_movement} negativo con usuario y guía.
    \item Si la despacho proviene de una solicitud, \texttt{SolicitudService} actualiza \texttt{cantidad\_despachada} y recalcula el estado (Pendiente/Parcial/Despachado).
    \item Se confirma la transacción y se muestra el correlativo. La guía es descargable en PDF.
\end{enumerate}

\subsection{Flujo alternativo}
\begin{itemize}
    \item \textbf{Stock insuficiente}: transacción revertida, mensaje con solicitado vs.\ disponible; la vista muestra el faltante.
    \item \textbf{Lote seleccionado sin stock suficiente}: el error identifica el lote y el estado del stock.
    \item \textbf{Producto controlado inyectado en el formulario}: rechazo explícito del servidor.
    \item \textbf{Solicitud con estupefacientes}: se excluyen del despacho normal y se informa que deben salir por el módulo controlado.
\end{itemize}

% =========================================================
\section{Reglas de negocio}
\begin{enumerate}
    \item \textbf{FEFO obligatorio server-side}: primero vence, primero sale; el cliente solo sugiere.
    \item \textbf{Estados de lote}: solo stock \texttt{Disponible} es despachable; \texttt{Cuarentena} y \texttt{Merma} quedan bloqueados y fuera de reportes de disponibilidad.
    \item \textbf{Segregación de salidas}: estupefacientes solo por Despacho Controlado (GDC) y solo para solicitudes CESFAM; el resto por Despacho Normal (GD).
    \item \textbf{Controlados}: lista fija de 16 fármacos/recetarios en \texttt{config/drogueria.php}, aplicada en recepción (categoría automática), despachos y solicitudes.
    \item \textbf{Acumulación de lotes}: una recepción o importación sobre un lote existente suma cantidad y actualiza valor/usuario; nunca duplica el registro.
    \item \textbf{IVA}: 19\% configurable (\texttt{config('drogueria.iva')}) sobre el valor neto.
    \item \textbf{Correlativos}: autogenerados con bloqueo y restricción UNIQUE (GD, GDC, recepción).
    \item \textbf{Trazabilidad inmutable}: todo ingreso/egreso registra usuario, lote, documento y cantidad con signo.
    \item \textbf{Establecimientos}: lista oficial de 15 centros; solo los 4 CESFAM solicitan estupefacientes.
    \item \textbf{Ocupación}: \texttt{ocupación\% = stock disponible del estante / capacidad de la estantería}, con semáforo $<$50\%, 50--79\%, $\geq$80\%.
\end{enumerate}

% =========================================================
\section{Diseño de API}

La aplicación expone endpoints JSON internos (mismo origen, autenticados por sesión):
\begin{table}[H]
\centering
\begin{tabularx}{\textwidth}{|l|l|X|}
\hline
\rowcolor{ucsblue} \textcolor{white}{\textbf{Endpoint}} & \textcolor{white}{\textbf{Método}} & \textcolor{white}{\textbf{Función}} \\ \hline
\texttt{/api/global-search?q=} & GET & Búsqueda global en inventario, catálogo y proveedores. \\ \hline
\texttt{/api/notifications} & GET & Alertas de vencimiento, cuarentena y stock bajo. \\ \hline
\texttt{/dispatch/solicitudes} & GET & Solicitudes pendientes/parciales en JSON para los despachos. \\ \hline
\texttt{/traceability/timeline?lote=} & GET & Línea de tiempo de un lote + ficha completa. \\ \hline
\texttt{/qr/\{inventory\}/image} & GET & Etiqueta QR en SVG del lote. \\ \hline
\texttt{/shelf/\{id\}/update-position} & POST & Persiste coordenadas 3D de una estantería. \\ \hline
\texttt{/suppliers/\{id\}/toggle-status} & PATCH & Activa/desactiva proveedor. \\ \hline
\end{tabularx}
\caption{Endpoints internos de la plataforma.}
\end{table}

% =========================================================
\section{Seguridad}
\begin{itemize}
    \item \textbf{Autenticación}: sesión con regeneración al ingresar, cookie cifrada, \texttt{throttle} en recuperación.
    \item \textbf{Autorización}: middleware \texttt{role:Administrador} en gestión de usuarios/roles; \texttt{permission:ver.<módulo>} en cada módulo; \texttt{Gate::before} para el Administrador; sidebar filtrado por \texttt{@can}.
    \item \textbf{Protección de formularios}: CSRF en el 100\% de los POST/PUT/DELETE, incluidas peticiones AJAX con \texttt{X-CSRF-TOKEN}.
    \item \textbf{Validación}: reglas server-side en todos los escritos (\texttt{in:}, \texttt{exists}, \texttt{min/max}, \texttt{Rule::in} para establecimientos).
    \item \textbf{Mass assignment}: \texttt{\$fillable} explícito en todos los modelos.
    \item \textbf{Inyección SQL}: Eloquent/Query Builder parametrizado; LIKEs escapados.
    \item \textbf{Datos sensibles}: contraseñas con bcrypt; errores genéricos al usuario con detalle técnico solo en \texttt{storage/logs}.
    \item \textbf{XSS}: salida escapada por Blade (dobles llaves); toasts con \texttt{@json}.
\end{itemize}

% =========================================================
\section{Estructura del proyecto}
\begin{lstlisting}
proyecto_drogueria/
|- app/
|  |- Http/Controllers/        # 18 controladores (modulos + Auth)
|  |- Imports/InventoriesImport.php
|  |- Exports/InventoryExport.php
|  |- Models/                  # 18 modelos Eloquent
|  |- Services/
|  |  |- InventoryService.php  # FEFO, ingresos, movimientos
|  |  |- SolicitudService.php  # despachos contra solicitudes
|- bootstrap/app.php           # alias middleware + handler permisos
|- config/drogueria.php        # IVA, controlados, establecimientos
|- database/
|  |- migrations/              # esquema versionado
|  |- seeders/                 # roles, usuarios, distribucion
|- public/plantillas/          # plantilla de carga masiva
|- resources/views/            # vistas Blade por modulo
|- routes/web.php              # rutas + permisos por modulo
\end{lstlisting}

% =========================================================
\section{Diseño de interfaces}

\subsection{Dashboard}
KPIs de valor de inventario, unidades, despachos y recepciones del período, próximos a vencer, cuarentena, últimos movimientos y panel QR del último lote; filtro de período (hoy/semana/mes/año).

\subsection{Recepción e Inventario}
Formularios por secciones numeradas, filas dinámicas de productos, totales automáticos, tablas agrupadas con desglose de lotes y modales operativos valorizados.

\subsection{Despachos y Distribución}
Formularios de guía con firmas digitales y escáner QR; visor 3D con leyenda de ocupación, panel de control y tooltips.

\subsection{Trazabilidad, Reportes y Usuarios}
Tablas con filtros avanzados, KPIs de valorización y modales de roles con matriz de accesos por vista.

% =========================================================
\section{Diagrama de secuencia}
Flujo de despacho con FEFO:

\begin{lstlisting}
Usuario        DespachoController    InventoryService       MySQL
   |-- POST /dispatch ----------->|                      |
   |                              |-- beginTransaction ->|
   |                              |-- generarCorrelativo (lockForUpdate)
   |                              |-- consumir() -------->|
   |                              |   SELECT lotes WHERE estado='Disponible'
   |                              |     AND producto=? FOR UPDATE ORDER BY
   |                              |     fecha_vencimiento ASC
   |                              |   UPDATE inventories SET cantidad=cantidad-N
   |                              |-- registrarMovimiento (stock_movements)
   |                              |-- SolicitudService::registrarDespacho()
   |                              |-- COMMIT ------------>|
   |<-- redirect + flash OK ------|
\end{lstlisting}

% =========================================================
\section{Diseño de despliegue}
\begin{itemize}
    \item \textbf{Entorno de desarrollo}: Laragon (Apache + PHP 8.3 + MySQL 8) en Windows; dominio local \texttt{proyecto\_drogueria.test}.
    \item \textbf{Producción sugerida}: servidor Linux con Nginx/Apache, PHP-FPM 8.3, MySQL 8, HTTPS obligatorio (requerido para la cámara QR fuera de localhost), \texttt{php artisan config:cache route:cache view:cache}, respaldo diario de la base de datos.
    \item \textbf{Variables}: \texttt{.env} con \texttt{APP\_KEY}, credenciales de BD y configuración de correo.
\end{itemize}

% =========================================================
\section{Control de versiones}
El historial completo de versiones del documento se encuentra en la sección ``Control de versiones'' al inicio. El código fuente se versiona en GitHub.

% =========================================================
\section{Plan de pruebas}
\begin{itemize}
    \item \textbf{Pruebas de servicio (13/13 OK)}: FEFO con reparto entre lotes, lote explícito, stock insuficiente, bloqueo de lotes en cuarentena, acumulación de lotes, movimientos con FKs, QR SVG.
    \item \textbf{Pruebas de importación (11/11 OK)}: parsing de fechas (texto y serial Excel), estados, acumulación, rechazo de filas inválidas con número de registro, enlaces a productos y usuario.
    \item \textbf{Pruebas HTTP end-to-end}: login y navegación de 16 vistas; recepción $\to$ liberación $\to$ despacho FEFO $\to$ trazabilidad; solicitud $\to$ despachos divididos normal/controlado $\to$ cierre de solicitud; QR; carga masiva por archivo real; permisos por rol (accesos permitidos y bloqueos por redirect).
    \item \textbf{Regresión visual}: verificación automatizada de 16 rutas con detección de errores de render.
\end{itemize}

% =========================================================
\section{Respaldo y recuperación}
\begin{itemize}
    \item Respaldo diario de la base de datos (\texttt{mysqldump}) con retención de 30 días.
    \item Respaldo semanal del directorio \texttt{storage/app} (firmas y archivos).
    \item Procedimiento de recuperación: restaurar el dump más reciente, verificar \texttt{.env} y \texttt{APP\_KEY}, ejecutar \texttt{php artisan migrate --force} y validar el módulo de inventario.
    \item La trazabilidad permite reconstruir el historial completo de un lote ante pérdida de datos parciales.
\end{itemize}

% =========================================================
\section{Mantenimiento}
\begin{itemize}
    \item Incorporar un estupefaciente nuevo: agregarlo a \texttt{config/drogueria.php} (lista \texttt{controlados}) --- recepción, despachos y solicitudes lo adoptan automáticamente.
    \item Incorporar un establecimiento: agregarlo a la lista \texttt{establecimientos} con su tipo (CESFAM/Otros) y dirección.
    \item Ajustar el IVA o los meses de alerta: \texttt{config/drogueria.php}.
    \item Nuevos permisos de vista: extender \texttt{UserController::MODULOS} y el sidebar con \texttt{@can}.
    \item Dependencias: \texttt{composer update} y \texttt{npm ci \&\& npm run build} con pruebas previas.
\end{itemize}

% =========================================================
\section{Conclusiones}
La plataforma cumple el objetivo de digitalizar integralmente la cadena de abastecimiento de la droguería comunal: recepción documentada, bodega valorizada por lotes con FEFO obligatorio, control sanitario (cuarentena/merma), segregación legal de estupefacientes, distribución física modelada en 3D con ocupación en tiempo real, trazabilidad completa por lote y reportes valorizados. La arquitectura en servicios y la configuración centralizada permiten evolucionar el sistema sin riesgo de duplicar reglas de negocio.

% =========================================================
\section{Anexos}

\subsection{Glosario}
\begin{itemize}
    \item \textbf{FEFO}: First Expired, First Out --- descuento del lote con fecha de vencimiento más próxima.
    \item \textbf{Lote}: unidad de trazabilidad física de un producto.
    \item \textbf{CESFAM}: Centro de Salud Familiar.
    \item \textbf{CECOSF}: Centro Comunitario de Salud Familiar.
    \item \textbf{SAR}: Servicio de Atención de Riesgo (SAR).
    \item \textbf{UAPO}: Unidad de Atención Primaria de Urgencia ( terminology local).
    \item \textbf{GD / GDC}: Guía de Despacho / Guía de Despacho Controlado.
    \item \textbf{FOFAR}: Fármacos de Oportunidad Farmacéutica.
    \item \textbf{Cuarentena}: retención sanitaria preventiva de un lote.
    \item \textbf{Merma}: baja definitiva de stock (daño, vencimiento).
    \item \textbf{ISP}: Instituto de Salud Pública de Chile.
\end{itemize}

\subsection{Configuración centralizada (config/drogueria.php)}
\begin{lstlisting}
'iva' => 0.19,
'meses_alerta_vencimiento' => 6,
'controlados' => [ /* 16 farmacos + recetarios */ ],
'establecimientos' => [
    ['nombre' => 'CESFAM Los Cerros',   'tipo' => 'CESFAM', 'direccion' => '...'],
    ['nombre' => 'CESFAM San Vicente',  'tipo' => 'CESFAM', 'direccion' => '...'],
    ['nombre' => 'CECOSF Centinela',    'tipo' => 'Otros',  'direccion' => '...'],
    // 15 establecimientos oficiales
],
\end{lstlisting}

\subsection{Comandos artisan relevantes}
\begin{lstlisting}
php artisan migrate --force                     # aplicar esquema
php artisan db:seed --class=DistributionSeeder  # bodega fisica
php artisan view:cache / route:cache            # optimizacion
php artisan optimize:clear                      # limpiar caches
\end{lstlisting}

\end{document}
